% Generated draft: requires mathematical and source review.
\documentclass[11pt]{article}
\usepackage[margin=25mm]{geometry}
\usepackage{fontspec}
\setmainfont{Latin Modern Roman}
\usepackage{amsmath,amssymb,mathtools,unicode-math}
\setmathfont{Latin Modern Math}
\usepackage{enumitem,needspace,xurl,hyperref}
\setlength{\parindent}{0pt}
\setlength{\parskip}{5pt}
\setlength{\emergencystretch}{3em}
\newcommand{\N}{\mathbb{N}}
\newcommand{\Z}{\mathbb{Z}}
\newcommand{\Q}{\mathbb{Q}}
\newcommand{\R}{\mathbb{R}}
\newcommand{\C}{\mathbb{C}}
\newcommand{\F}{\mathbb{F}}
\newcommand{\A}{\mathbb{A}}
\newcommand{\PP}{\mathbb{P}}
\newcommand{\E}{\mathbb{E}}
\newcommand{\eps}{\varepsilon}
\newcommand{\dd}{\,\mathrm{d}}
\begin{document}
\section*{UnsolvedMath ID: 2307070}
\textbf{Source problem code:} catalogue entry
\subsection{Definitions and mathematical statement}
The source uses a Clifford-algebra model of M\"obius transformations of $\R^n$, writing
\[
\phi(x)=(ax+b)(cx+d)^{-1},
\]
where $a,b,c,d\in A_n$ obey constraints not reproduced in the record. The inverse and all expressions are restricted to points where defined. In a specified Euclidean Clifford convention, the Dirac operator is $D=\sum_{j=1}^n e_j\partial_{x_j}$; the precise algebra, vector embedding, and generator signs must be fixed in that convention.

The source reports a linear conformally invariant class in the printed form
\[
D^k f_k\bigl((ax+b)(cx+d)^{-1}\bigr)=0,\qquad k\in\N,
\]
with associated weights
\[
J_k(cx+d)=
\begin{cases}
(cx+d)^*|cx+d|^{-n-1+k},&k=2p-1,\\
|cx+d|^{-n+k},&k=2p.
\end{cases}
\]
Here $*$ denotes a Clifford involution referred to by the source but not defined; it is not identified silently with conjugation or reversion. The coefficient constraints and norm convention are also absent. The record does not explain how the associated weight is inserted into the printed transformed equation, so the formula is retained as supplied rather than promoted to a fully specified covariance identity.

J. Ryan asks:
\begin{enumerate}
\item Which nonlinear differential equations have conformally invariant solution spaces?
\item Can their conformal weights likewise be expressed in terms of $cx+d$, and what relations do their solutions have to those of the linear conformally invariant equations?
\end{enumerate}
In a precise formulation, conformal invariance would require a specified transformation law, often a weighted pullback such as $u\mapsto J_\phi(u\circ\phi)$, to carry solutions on the target domain to solutions on the source domain. In the noncommutative algebra the multiplication order is part of that law. The problem does not fix the admissible nonlinearities, differential orders, dependent-variable type, regularity, or meaning of classification and relation to linear solutions. Thus it is a source-specific classification program, not one universal theorem with supplied hypotheses. The summary reports it open as of 2018, with no definitive resolution located.
\subsection{Short English statement}
Ryan asks for nonlinear analogues of Clifford-Dirac equations whose solution spaces transform conformally, together with their weights and relation to linear solutions. The source omits the Clifford constraints, involution, and precise nonlinear class and transformation law, so these cannot be invented from the linear weight formula.
\subsection{Sources}
\begin{itemize}
\item Local UnsolvedMath registry, permanent ID 2307070, exact source code catalogue entry, ``Research Problems in Function Theory --- Problem 7.70''; status open, classification OPEN-TRIAGE.
\item Hayman, \emph{Research Problems in Function Theory} (2018), Problem 7.70, attributed to J. Ryan, \url{https://arxiv.org/abs/1809.07200}. The record identifies extraction from public arXiv source TeX, accessed July 29, 2026; the 2018 update was retained for status triage.
\item Ahlfors is credited for the M\"obius representation, but the parenthetical reference is empty. The involution and reported linear classification likewise refer to missing source material, with no exact theorem or full citation supplied.
\item The supplied research summary reports openness as of the 2018 edition and no definitive resolution located, with no additional authoritative reference or precise search date.
\end{itemize}
\end{document}
